% Compile with: pdflatex lenses_formulas.tex
% One math card per lens — presentation-ready, screenshot directly into PPT.

\documentclass[12pt]{article}
\usepackage[a4paper, margin=18mm]{geometry}
\usepackage{amsmath, amssymb}
\usepackage{xcolor}
\usepackage{tcolorbox}
\tcbuselibrary{skins, breakable}
\usepackage{parskip}
\pagestyle{empty}

\definecolor{nlf}{RGB}{240, 180, 100}
\definecolor{cfa}{RGB}{120, 180, 110}
\definecolor{res}{RGB}{100, 150, 220}
\definecolor{lca}{RGB}{220, 120, 150}

\newtcolorbox{lensbox}[2]{%
  enhanced, breakable,
  colback=#2!8, colframe=#2!80!black,
  fonttitle=\bfseries\large,
  title={#1},
  arc=2mm, boxrule=0.6pt,
  left=4mm, right=4mm, top=2mm, bottom=2mm,
}

\newcommand{\say}[1]{\par\smallskip\textit{\small \textcolor{black!60}{What to say:} #1}}

\begin{document}

\begin{center}
{\Huge\bfseries Forensic lenses --- math cheat sheet}\\[1mm]
{\small ADAIS / wave\_1 ~~~~ one slide per lens}
\end{center}
\bigskip

% =========================================================================
\begin{lensbox}{NLF --- Noise-Level Function}{nlf}

For every $8{\times}8$ patch $p$ of the gray image, compute mean and std:
\[
\bar I_p \;=\; \frac{1}{64}\sum_{(i,j)\in p} I(i,j),
\qquad
\sigma_p \;=\; \sqrt{\frac{1}{64}\sum_{(i,j)\in p}\bigl(I(i,j)-\bar I_p\bigr)^{\!2}}.
\]
Throw away the most textured 50\% of patches (high local gradient).
Bin the rest by $\bar I_p$ into 16 intensity bins; for each bin take the
median $\sigma$. Fit the \textbf{Poisson--Gaussian sensor model}:
\[
\boxed{\;\sigma^{2}(I)\;=\;a\cdot I \;+\; b\;}
\qquad
R^{2} \;=\; 1 \;-\; \frac{\sum_i (\sigma_i^{2}-\hat\sigma_i^{2})^{2}}
                          {\sum_i (\sigma_i^{2}-\overline{\sigma^{2}})^{2}}.
\]
\textbf{Output (20 values):} the 16 bin medians (the curve) plus
$(a,\,b,\,R^{2},\,\text{valid\_frac})$.

\say{Real cameras have \emph{shot noise} (Poisson) that grows with
brightness, plus a constant \emph{read noise} (Gaussian). The line should
fit. AI generators don't simulate sensor physics, so the curve is flat or
noisy and $R^{2}$ is low.}
\end{lensbox}

\bigskip
% =========================================================================
\begin{lensbox}{CFA --- Colour-Filter-Array periodicity}{cfa}

For each channel $c\in\{R,G,B\}$, take the high-pass residual
(subtract a Gaussian blur with $\sigma{=}1$):
\[
r_c \;=\; c \;-\; G_{\sigma}\!\ast\, c.
\]
Take the 2-D FFT magnitude (shifted so DC is at the centre):
\[
S_c(\omega_y,\omega_x) \;=\; \bigl|\mathcal{F}\{r_c\}\bigr|.
\]
Sample three diagnostic peaks at the half-Nyquist frequencies and
normalise by the mean spectrum so amplitude doesn't matter:
\[
\boxed{\;
\tilde p_{c,k} \;=\;
\frac{\max\limits_{\|(\omega_y,\omega_x)-q_k\|\le 2}\!\!S_c(\omega_y,\omega_x)}
     {\overline{S_c}}\;},
\quad
q_k \in \Bigl\{\bigl(0,\tfrac{\pi}{2}\bigr),\bigl(\tfrac{\pi}{2},0\bigr),
              \bigl(\tfrac{\pi}{2},\tfrac{\pi}{2}\bigr)\Bigr\}.
\]
\textbf{Output (12 values):} 3 channels $\times$ \{hh, vh, dh, pool\}.

\say{Cameras only sense one colour per pixel (Bayer mosaic) and
\emph{interpolate} the missing two from neighbours. That interpolation
leaves a 2-pixel periodicity, which shows up as a spike at $\pi/2$ in
the FFT. Generators emit RGB directly --- no spike.}
\end{lensbox}

\bigskip
% =========================================================================
\begin{lensbox}{RES (radial) --- residual spectrum, radial bins}{res}

Take the denoising residual of the gray image
$\;r = I - D_{\text{wavelet}}(I)$, then its FFT magnitude:
\[
S(\omega_y,\omega_x) \;=\; \bigl|\mathcal{F}\{r\}\bigr|,
\qquad
\rho(\omega_y,\omega_x) \;=\; \sqrt{(\omega_y-c_y)^{2}+(\omega_x-c_x)^{2}}.
\]
Define 64 \emph{log-spaced} radius bins
$1 = \rho_0 < \rho_1 < \dots < \rho_{64} = N/2-1$. The feature for bin $k$
is the mean spectrum on that ring, L1-normalised:
\[
\boxed{\;
P_k \;=\; \frac{1}{|B_k|}\!\!\sum_{(\omega_y,\omega_x)\in B_k}\!\!
            S(\omega_y,\omega_x)\;},
\qquad
\tilde P_k = \frac{P_k}{\sum_j P_j},
\]
\[
B_k = \{(\omega_y,\omega_x):\rho_k\le \rho(\omega_y,\omega_x)<\rho_{k+1}\}.
\]
\textbf{Output:} 64 values --- the \emph{shape} of the radial spectrum.

\say{Real photos have a smooth $1/f$ fall-off. GAN and diffusion
\emph{upsamplers} leave periodic copies of the spectrum --- visible as
bumps at specific radii. The radial profile is the cleanest way to spot
them.}
\end{lensbox}

\clearpage
% =========================================================================
\begin{lensbox}{RES (angular) --- residual spectrum, angular bins}{res}

Same residual $r$ and spectrum $S$ as above, but now bin by orientation
inside a mid-frequency annulus:
\[
\theta(\omega_y,\omega_x) \;=\; \operatorname{atan2}(\omega_y-c_y,\;\omega_x-c_x)\;\bmod\;\pi
\quad\text{(fold to }[0,\pi)\text{)}.
\]
Mask to the annulus $\mathcal{A} = \{0.25\,r_{\text{nyq}} \le \rho \le 0.75\,r_{\text{nyq}}\}$
and bin into 16 sectors of $11.25^{\circ}$:
\[
\boxed{\;
Q_k \;=\; \frac{1}{|\mathcal{A}_k|}\!\!\sum_{(\omega_y,\omega_x)\in\mathcal{A}_k}
   \!\! S(\omega_y,\omega_x)\;},
\qquad
\tilde Q_k = \frac{Q_k}{\sum_j Q_j},
\]
\[
\mathcal{A}_k = \mathcal{A}\,\cap\,\{\theta\in[\theta_k,\theta_{k+1})\}.
\]
\textbf{Output:} 16 values --- the orientation pattern of the residual.

\say{A real photo is roughly \emph{isotropic} in the mid-frequencies ---
all directions look the same. Each generator's upsampling kernel has a
preferred axis (horizontal, vertical, $45^{\circ}$ lobes), so the angular
profile is what tells \emph{LlamaGen vs.\ RAR vs.\ VAR} apart.}
\end{lensbox}

\bigskip
% =========================================================================
\begin{lensbox}{LCA --- Lateral Chromatic Aberration}{lca}

Split the image into a $4{\times}4$ grid of tiles. For each tile, align
the R and B channels to G by integer cross-correlation (within $\pm 3$ px):
\[
(\delta y,\delta x)^{*}
\;=\;\operatorname*{argmax}_{|\delta y|,|\delta x|\le 3}
\frac{\sum (c_{\,\delta}-\bar c)(G-\bar G)}{\|c_{\,\delta}-\bar c\|\,\|G-\bar G\|}
\quad (\text{normalised cross-correlation}).
\]
Each tile centre has position $(y_t,x_t)$ relative to the image centre,
distance $d_t=\sqrt{y_t^{2}+x_t^{2}}$, outward unit vector $\hat n_t$.
Decompose the shift into radial / tangential components:
\[
\Delta_{\text{rad},t} \;=\; (\delta y,\delta x)\!\cdot\!\hat n_t,
\qquad
\Delta_{\text{tan},t} \;=\; (\delta y,\delta x)\!\cdot\!\hat n_t^{\perp}.
\]
Fit a line through the origin (LCA is zero at the optical centre):
\[
\boxed{\;
\hat s \;=\; \frac{\sum_t \Delta_{\text{rad},t}\,d_t}{\sum_t d_t^{\,2}}\;},
\qquad
R^{2} \;=\; 1-\frac{\sum_t(\Delta_{\text{rad},t}-\hat s\, d_t)^{2}}
                   {\sum_t(\Delta_{\text{rad},t}-\overline{\Delta_{\text{rad}}})^{2}}.
\]
\textbf{Output (8 values):} two pairs (R$\to$G, B$\to$G) $\times$
$(\hat s,\, R^{2},\, \overline{|\Delta_{\text{tan}}|},\, \operatorname{std}(\Delta_{\text{rad}}))$.

\say{Glass refracts red and blue light by different angles, so a real
lens shifts the R and B channels outward from the optical centre, with
the shift growing \emph{linearly} with distance. Real $\to$ high slope,
high $R^{2}$, near-zero tangential. AI $\to$ random scatter, $R^{2}\!\approx\!0$.}
\end{lensbox}

\bigskip

\begin{center}
\fbox{\parbox{0.92\textwidth}{\centering
\textbf{The one-sentence summary for the audience:}\\[2pt]
\itshape ``Every formula here measures something a real camera \emph{has}
that a generative model \emph{doesn't simulate}: sensor noise statistics,
demosaicing periodicity, upsampling-free spectrum, and lens-glass
chromatic dispersion.''
}}
\end{center}

\end{document}
